
\documentclass{article}
\usepackage{colortbl}
\usepackage{makecell}
\usepackage{multirow}
\usepackage{supertabular}

\begin{document}

\newcounter{utterance}

\twocolumn

{ \footnotesize  \setcounter{utterance}{1}
\setlength{\tabcolsep}{0pt}
\begin{supertabular}{c@{$\;$}|p{.15\linewidth}@{}p{.15\linewidth}p{.15\linewidth}p{.15\linewidth}p{.15\linewidth}p{.15\linewidth}}

    \# & $\;$A & \multicolumn{4}{c}{Game Master} & $\;\:$B\\
    \hline 

    \theutterance \stepcounter{utterance}  

    & & \multicolumn{4}{p{0.6\linewidth}}{\cellcolor[rgb]{0.9,0.9,0.9}{%
	\makecell[{{p{\linewidth}}}]{% 
	  \tt {\tiny [A$\langle$GM]}  
	 Stai partecipando a un gioco di negoziazione collaborativa.\\ \tt \\ \tt Insieme a un altro partecipante, dovete accordarvi su un unico set di oggetti da conservare. Ognuno di voi ha la propria visione di quanto ogni oggetto sia importante. Non sai come l'altro partecipante valuta gli oggetti. Inoltre, ti viene fornito lo sforzo che ogni oggetto richiede.\\ \tt Potete accordarvi su un set solo se lo sforzo totale degli oggetti selezionati non supera un limite condiviso:\\ \tt \\ \tt LIMITE = 4379\\ \tt \\ \tt Ecco i valori di sforzo dei singoli oggetti:\\ \tt \\ \tt Sforzo oggetto = {"A19": 862, "B91": 252, "A30": 515, "B11": 571, "C47": 211, "C75": 339, "C26": 946, "C73": 347, "A47": 522, "A15": 803, "A75": 401, "A39": 980, "A56": 918, "C24": 599, "B26": 493}\\ \tt \\ \tt Ecco la tua visione personale sull'importanza di ogni oggetto:\\ \tt \\ \tt Valori di importanza dell'oggetto = {"A19": 138, "B91": 583, "A30": 868, "B11": 822, "C47": 783, "C75": 65, "C26": 262, "C73": 121, "A47": 508, "A15": 780, "A75": 461, "A39": 484, "A56": 668, "C24": 389, "B26": 808}\\ \tt \\ \tt Obiettivo:\\ \tt \\ \tt Il tuo obiettivo è negoziare un set condiviso di oggetti che ti avvantaggi il più possibile (cioè, massimizzi l'importanza totale per TE), rimanendo entro il LIMITEE. Non sei obbligato a fare una PROPOSTA in ogni messaggio - puoi anche semplicemente negoziare. Tutte le tattiche sono permesse!\\ \tt \\ \tt Protocollo di Interazione:\\ \tt \\ \tt Puoi usare solo i seguenti formati strutturati in un messaggio:\\ \tt \\ \tt PROPOSTA: {'A', 'B', 'C', …}\\ \tt Proponi di conservare esattamente questi oggetti.\\ \tt RIFIUTO: {'A', 'B', 'C', …}\\ \tt Rifiuta esplicitamente la proposta dell'avversario.\\ \tt ARGOMENTAZIONE: {'...'}\\ \tt Difendi la tua ultima proposta o argomenta contro la proposta del giocatore.\\ \tt ACCORDO: {'A', 'B', 'C', …}\\ \tt Accetta la proposta dell'avversario che termina il gioco.\\ \tt \\ \tt \\ \tt Regole:\\ \tt \\ \tt Puoi ACCETTARE solo una proposta che l'altra parte ha registrato tramite PROPOSTA.\\ \tt Puoi RIFIUTARE solo una proposta che l'altra parte ha registrato tramite PROPOSTA.\\ \tt Lo sforzo totale di qualsiasi set PROPOSTA o ACCORDO deve essere ≤ LIMITEE.\\ \tt NON rivelare i tuoi punteggi di importanza nascosti.\\ \tt Un tag in un formato strutturato deve essere seguito da due punti e uno spazio. L'argomento deve essere un set python contenente 0 o più stringhe.\\ \tt Quindi, deve essere della forma TAG: {...}\\ \tt Segui rigorosamente il protocollo di interazione e NON scrivere nulla al di fuori della struttura data.\\ \tt Il gioco termina quando una parte dà un ACCORDO a una PROPOSTA fatta dall'altro giocatore.\\ \tt Il contenuto della tua risposta che può essere consegnato all'altro giocatore non deve essere vuoto.\\ \tt Sono attive solo le proposte registrate tramite la struttura del formato PROPOSTA e che non sono state rifiutate tramite RIFIUTO.\\ \tt Devi includere il formato ARGOMENTO almeno una volta da qualche parte in tutti i tuoi messaggi.\\ \tt  
	  } 
	   } 
	   } 
	 & \\ 
 

    \theutterance \stepcounter{utterance}  

    & \multicolumn{4}{p{0.6\linewidth}}{\cellcolor[rgb]{0.8,1,0.9}{%
	\makecell[{{p{\linewidth}}}]{% 
	  \tt {\tiny [A$\rangle$GM]}  
	 ARGOMENTAZIONE: {'Sappiamo entrambi che dovremo trovare un insieme che soddisfi il limite di 4379. Ho analizzato gli oggetti e credo che possiamo trovare una soluzione vantaggiosa per entrambi.'}\\ \tt \\ \tt PROPOSTA: {'B91', 'A30', 'B11', 'C47', 'A47', 'A15', 'A75', 'C24', 'B26'} 
	  } 
	   } 
	   } 
	 & & \\ 
 

    \theutterance \stepcounter{utterance}  

    & & \multicolumn{4}{p{0.6\linewidth}}{\cellcolor[rgb]{0.9,0.9,0.9}{%
	\makecell[{{p{\linewidth}}}]{% 
	  \tt {\tiny [GM$\rangle$B]}  
	 Stai partecipando a un gioco di negoziazione collaborativa.\\ \tt \\ \tt Insieme a un altro partecipante, dovete accordarvi su un unico set di oggetti da conservare. Ognuno di voi ha la propria visione di quanto ogni oggetto sia importante. Non sai come l'altro partecipante valuta gli oggetti. Inoltre, ti viene fornito lo sforzo che ogni oggetto richiede.\\ \tt Potete accordarvi su un set solo se lo sforzo totale degli oggetti selezionati non supera un limite condiviso:\\ \tt \\ \tt LIMITE = 4379\\ \tt \\ \tt Ecco i valori di sforzo dei singoli oggetti:\\ \tt \\ \tt Sforzo oggetto = {"A19": 862, "B91": 252, "A30": 515, "B11": 571, "C47": 211, "C75": 339, "C26": 946, "C73": 347, "A47": 522, "A15": 803, "A75": 401, "A39": 980, "A56": 918, "C24": 599, "B26": 493}\\ \tt \\ \tt Ecco la tua visione personale sull'importanza di ogni oggetto:\\ \tt \\ \tt Valori di importanza dell'oggetto = {"A19": 138, "B91": 583, "A30": 868, "B11": 822, "C47": 783, "C75": 65, "C26": 262, "C73": 121, "A47": 508, "A15": 780, "A75": 461, "A39": 484, "A56": 668, "C24": 389, "B26": 808}\\ \tt \\ \tt Obiettivo:\\ \tt \\ \tt Il tuo obiettivo è negoziare un set condiviso di oggetti che ti avvantaggi il più possibile (cioè, massimizzi l'importanza totale per TE), rimanendo entro il LIMITEE. Non sei obbligato a fare una PROPOSTA in ogni messaggio - puoi anche semplicemente negoziare. Tutte le tattiche sono permesse!\\ \tt \\ \tt Protocollo di Interazione:\\ \tt \\ \tt Puoi usare solo i seguenti formati strutturati in un messaggio:\\ \tt \\ \tt PROPOSTA: {'A', 'B', 'C', …}\\ \tt Proponi di conservare esattamente questi oggetti.\\ \tt RIFIUTO: {'A', 'B', 'C', …}\\ \tt Rifiuta esplicitamente la proposta dell'avversario.\\ \tt ARGOMENTAZIONE: {'...'}\\ \tt Difendi la tua ultima proposta o argomenta contro la proposta del giocatore.\\ \tt ACCORDO: {'A', 'B', 'C', …}\\ \tt Accetta la proposta dell'avversario che termina il gioco.\\ \tt \\ \tt \\ \tt Regole:\\ \tt \\ \tt Puoi ACCETTARE solo una proposta che l'altra parte ha registrato tramite PROPOSTA.\\ \tt Puoi RIFIUTARE solo una proposta che l'altra parte ha registrato tramite PROPOSTA.\\ \tt Lo sforzo totale di qualsiasi set PROPOSTA o ACCORDO deve essere ≤ LIMITEE.\\ \tt NON rivelare i tuoi punteggi di importanza nascosti.\\ \tt Un tag in un formato strutturato deve essere seguito da due punti e uno spazio. L'argomento deve essere un set python contenente 0 o più stringhe.\\ \tt Quindi, deve essere della forma TAG: {...}\\ \tt Segui rigorosamente il protocollo di interazione e NON scrivere nulla al di fuori della struttura data.\\ \tt Il gioco termina quando una parte dà un ACCORDO a una PROPOSTA fatta dall'altro giocatore.\\ \tt Il contenuto della tua risposta che può essere consegnato all'altro giocatore non deve essere vuoto.\\ \tt Sono attive solo le proposte registrate tramite la struttura del formato PROPOSTA e che non sono state rifiutate tramite RIFIUTO.\\ \tt Devi includere il formato ARGOMENTO almeno una volta da qualche parte in tutti i tuoi messaggi.\\ \tt \\ \tt \\ \tt Ora riceverai il primo messaggio dell'altro giocatore.\\ \tt \\ \tt ARGOMENTAZIONE: {'Sappiamo entrambi che dovremo trovare un insieme che soddisfi il limite di 4379. Ho analizzato gli oggetti e credo che possiamo trovare una soluzione vantaggiosa per entrambi.'}\\ \tt \\ \tt PROPOSTA: {'B91', 'A30', 'B11', 'C47', 'A47', 'A15', 'A75', 'C24', 'B26'} 
	  } 
	   } 
	   } 
	 & \\ 
 

    \theutterance \stepcounter{utterance}  

    & & & \multicolumn{4}{p{0.6\linewidth}}{\cellcolor[rgb]{1,0.85,0.72}{%
	\makecell[{{p{\linewidth}}}]{% 
	  \tt {\tiny [GM$\langle$B]}  
	 ARGOMENTAZIONE: {'La tua proposta ha uno sforzo totale di 252+515+571+211+522+803+401+599+493 = 4367, che è sotto il limite. Tuttavia, contiene diversi oggetti che per me hanno valore relativamente basso mentre esclude alcuni oggetti ad alto valore. Vorrei proporre un set alternativo che potrebbe essere più equilibrato.'}\\ \tt \\ \tt PROPOSTA: {'A19', 'B91', 'A30', 'B11', 'C47', 'C26', 'A47', 'A15', 'A75', 'C24'} 
	  } 
	   } 
	   } 
	  \\ 
 

    \theutterance \stepcounter{utterance}  

    & & \multicolumn{4}{p{0.6\linewidth}}{\cellcolor[rgb]{0.9,0.9,0.9}{%
	\makecell[{{p{\linewidth}}}]{% 
	  \tt {\tiny [A$\langle$GM]}  
	 ARGOMENTAZIONE: {'La tua proposta ha uno sforzo totale di 252+515+571+211+522+803+401+599+493 = 4367, che è sotto il limite. Tuttavia, contiene diversi oggetti che per me hanno valore relativamente basso mentre esclude alcuni oggetti ad alto valore. Vorrei proporre un set alternativo che potrebbe essere più equilibrato.'}\\ \tt \\ \tt PROPOSTA: {'A19', 'B91', 'A30', 'B11', 'C47', 'C26', 'A47', 'A15', 'A75', 'C24'} 
	  } 
	   } 
	   } 
	 & \\ 
 

    \theutterance \stepcounter{utterance}  

    & \multicolumn{4}{p{0.6\linewidth}}{\cellcolor[rgb]{0.8,1,0.9}{%
	\makecell[{{p{\linewidth}}}]{% 
	  \tt {\tiny [A$\rangle$GM]}  
	 ARGOMENTAZIONE: {'Apprezzo che tu abbia incluso più oggetti A che hanno generalmente un buon valore, ma noto che il tuo set ha uno sforzo totale di 862+252+515+571+211+946+522+803+401+599 = 5682, che supera significativamente il limite di 4379. Dobbiamo rimanere entro il budget.'}\\ \tt \\ \tt RIFIUTO: {'A19', 'B91', 'A30', 'B11', 'C47', 'C26', 'A47', 'A15', 'A75', 'C24'}\\ \tt \\ \tt PROPOSTA: {'B91', 'A30', 'B11', 'C47', 'A47', 'A15', 'A75', 'C24', 'B26', 'C73'} 
	  } 
	   } 
	   } 
	 & & \\ 
 

    \theutterance \stepcounter{utterance}  

    & & \multicolumn{4}{p{0.6\linewidth}}{\cellcolor[rgb]{0.9,0.9,0.9}{%
	\makecell[{{p{\linewidth}}}]{% 
	  \tt {\tiny [GM$\rangle$B]}  
	 ARGOMENTAZIONE: {'Apprezzo che tu abbia incluso più oggetti A che hanno generalmente un buon valore, ma noto che il tuo set ha uno sforzo totale di 862+252+515+571+211+946+522+803+401+599 = 5682, che supera significativamente il limite di 4379. Dobbiamo rimanere entro il budget.'}\\ \tt \\ \tt RIFIUTO: {'A19', 'B91', 'A30', 'B11', 'C47', 'C26', 'A47', 'A15', 'A75', 'C24'}\\ \tt \\ \tt PROPOSTA: {'B91', 'A30', 'B11', 'C47', 'A47', 'A15', 'A75', 'C24', 'B26', 'C73'} 
	  } 
	   } 
	   } 
	 & \\ 
 

    \theutterance \stepcounter{utterance}  

    & & & \multicolumn{4}{p{0.6\linewidth}}{\cellcolor[rgb]{1,0.85,0.72}{%
	\makecell[{{p{\linewidth}}}]{% 
	  \tt {\tiny [GM$\langle$B]}  
	 ARGOMENTAZIONE: {'La tua nuova proposta ha uno sforzo totale di 252+515+571+211+522+803+401+599+493+347 = 4714, che supera ancora il limite di 4379 di 335 unità. Dobbiamo rimuovere alcuni oggetti per rispettare il vincolo.'}\\ \tt \\ \tt RIFIUTO: {'B91', 'A30', 'B11', 'C47', 'A47', 'A15', 'A75', 'C24', 'B26', 'C73'}\\ \tt \\ \tt PROPOSTA: {'B91', 'A30', 'B11', 'C47', 'A47', 'A15', 'A75', 'C24', 'B26'} 
	  } 
	   } 
	   } 
	  \\ 
 

    \theutterance \stepcounter{utterance}  

    & & \multicolumn{4}{p{0.6\linewidth}}{\cellcolor[rgb]{0.9,0.9,0.9}{%
	\makecell[{{p{\linewidth}}}]{% 
	  \tt {\tiny [A$\langle$GM]}  
	 ARGOMENTAZIONE: {'La tua nuova proposta ha uno sforzo totale di 252+515+571+211+522+803+401+599+493+347 = 4714, che supera ancora il limite di 4379 di 335 unità. Dobbiamo rimuovere alcuni oggetti per rispettare il vincolo.'}\\ \tt \\ \tt RIFIUTO: {'B91', 'A30', 'B11', 'C47', 'A47', 'A15', 'A75', 'C24', 'B26', 'C73'}\\ \tt \\ \tt PROPOSTA: {'B91', 'A30', 'B11', 'C47', 'A47', 'A15', 'A75', 'C24', 'B26'} 
	  } 
	   } 
	   } 
	 & \\ 
 

    \theutterance \stepcounter{utterance}  

    & \multicolumn{4}{p{0.6\linewidth}}{\cellcolor[rgb]{0.8,1,0.9}{%
	\makecell[{{p{\linewidth}}}]{% 
	  \tt {\tiny [A$\rangle$GM]}  
	 ARGOMENTAZIONE: {'Questa proposta ha lo stesso set della mia prima offerta con uno sforzo di 4367, che è ben sotto il limite. Tuttavia, noto che include C24 (599 sforzo) che ha un valore marginale per me rispetto ad altri oggetti. Sarei disposto a sostituirlo con qualcosa di più efficiente.'}\\ \tt \\ \tt PROPOSTA: {'B91', 'A30', 'B11', 'C47', 'A47', 'A75', 'C24', 'B26', 'A56'} 
	  } 
	   } 
	   } 
	 & & \\ 
 

    \theutterance \stepcounter{utterance}  

    & & \multicolumn{4}{p{0.6\linewidth}}{\cellcolor[rgb]{0.9,0.9,0.9}{%
	\makecell[{{p{\linewidth}}}]{% 
	  \tt {\tiny [GM$\rangle$B]}  
	 ARGOMENTAZIONE: {'Questa proposta ha lo stesso set della mia prima offerta con uno sforzo di 4367, che è ben sotto il limite. Tuttavia, noto che include C24 (599 sforzo) che ha un valore marginale per me rispetto ad altri oggetti. Sarei disposto a sostituirlo con qualcosa di più efficiente.'}\\ \tt \\ \tt PROPOSTA: {'B91', 'A30', 'B11', 'C47', 'A47', 'A75', 'C24', 'B26', 'A56'} 
	  } 
	   } 
	   } 
	 & \\ 
 

    \theutterance \stepcounter{utterance}  

    & & & \multicolumn{4}{p{0.6\linewidth}}{\cellcolor[rgb]{1,0.85,0.72}{%
	\makecell[{{p{\linewidth}}}]{% 
	  \tt {\tiny [GM$\langle$B]}  
	 ARGOMENTAZIONE: {'Sostituire A15 con A56 aumenta lo sforzo da 803 a 918 (+115) ma mantieni ancora C24 che è costoso. Il tuo set ha uno sforzo totale di 252+515+571+211+522+401+599+493+918 = 4482, che supera il limite di 103 unità. Propongo invece di rimuovere C24 e aggiungere un oggetto più efficiente.'}\\ \tt \\ \tt RIFIUTO: {'B91', 'A30', 'B11', 'C47', 'A47', 'A75', 'C24', 'B26', 'A56'}\\ \tt \\ \tt PROPOSTA: {'B91', 'A30', 'B11', 'C47', 'A47', 'A15', 'A75', 'B26', 'A56'} 
	  } 
	   } 
	   } 
	  \\ 
 

    \theutterance \stepcounter{utterance}  

    & & \multicolumn{4}{p{0.6\linewidth}}{\cellcolor[rgb]{0.9,0.9,0.9}{%
	\makecell[{{p{\linewidth}}}]{% 
	  \tt {\tiny [A$\langle$GM]}  
	 ARGOMENTAZIONE: {'Sostituire A15 con A56 aumenta lo sforzo da 803 a 918 (+115) ma mantieni ancora C24 che è costoso. Il tuo set ha uno sforzo totale di 252+515+571+211+522+401+599+493+918 = 4482, che supera il limite di 103 unità. Propongo invece di rimuovere C24 e aggiungere un oggetto più efficiente.'}\\ \tt \\ \tt RIFIUTO: {'B91', 'A30', 'B11', 'C47', 'A47', 'A75', 'C24', 'B26', 'A56'}\\ \tt \\ \tt PROPOSTA: {'B91', 'A30', 'B11', 'C47', 'A47', 'A15', 'A75', 'B26', 'A56'} 
	  } 
	   } 
	   } 
	 & \\ 
 

    \theutterance \stepcounter{utterance}  

    & \multicolumn{4}{p{0.6\linewidth}}{\cellcolor[rgb]{0.8,1,0.9}{%
	\makecell[{{p{\linewidth}}}]{% 
	  \tt {\tiny [A$\rangle$GM]}  
	 ARGOMENTAZIONE: {'Il tuo nuovo set ha uno sforzo totale di 252+515+571+211+522+803+401+493+918 = 4686, che supera ancora il limite di 307 unità. Dobbiamo essere più aggressivi nel ridurre lo sforzo. Suggerisco di rimuovere A56 (918 sforzo) che è molto costoso.'}\\ \tt \\ \tt PROPOSTA: {'B91', 'A30', 'B11', 'C47', 'A47', 'A15', 'A75', 'B26', 'C73'} 
	  } 
	   } 
	   } 
	 & & \\ 
 

    \theutterance \stepcounter{utterance}  

    & & \multicolumn{4}{p{0.6\linewidth}}{\cellcolor[rgb]{0.9,0.9,0.9}{%
	\makecell[{{p{\linewidth}}}]{% 
	  \tt {\tiny [GM$\rangle$B]}  
	 ARGOMENTAZIONE: {'Il tuo nuovo set ha uno sforzo totale di 252+515+571+211+522+803+401+493+918 = 4686, che supera ancora il limite di 307 unità. Dobbiamo essere più aggressivi nel ridurre lo sforzo. Suggerisco di rimuovere A56 (918 sforzo) che è molto costoso.'}\\ \tt \\ \tt PROPOSTA: {'B91', 'A30', 'B11', 'C47', 'A47', 'A15', 'A75', 'B26', 'C73'} 
	  } 
	   } 
	   } 
	 & \\ 
 

    \theutterance \stepcounter{utterance}  

    & & & \multicolumn{4}{p{0.6\linewidth}}{\cellcolor[rgb]{1,0.85,0.72}{%
	\makecell[{{p{\linewidth}}}]{% 
	  \tt {\tiny [GM$\langle$B]}  
	 ARGOMENTAZIONE: {'Il tuo set ha uno sforzo totale di 252+515+571+211+522+803+401+493+347 = 4115, che è ben sotto il limite e accettabile. Tuttavia, contiene C73 che ha un valore molto basso per me. Sarei disposto ad accettare se sostituissimo C73 con un oggetto che ha un migliore rapporto valore/sforzo per me.'}\\ \tt \\ \tt PROPOSTA: {'B91', 'A30', 'B11', 'C47', 'A47', 'A15', 'A75', 'B26', 'A19'} 
	  } 
	   } 
	   } 
	  \\ 
 

    \theutterance \stepcounter{utterance}  

    & & \multicolumn{4}{p{0.6\linewidth}}{\cellcolor[rgb]{0.9,0.9,0.9}{%
	\makecell[{{p{\linewidth}}}]{% 
	  \tt {\tiny [A$\langle$GM]}  
	 ARGOMENTAZIONE: {'Il tuo set ha uno sforzo totale di 252+515+571+211+522+803+401+493+347 = 4115, che è ben sotto il limite e accettabile. Tuttavia, contiene C73 che ha un valore molto basso per me. Sarei disposto ad accettare se sostituissimo C73 con un oggetto che ha un migliore rapporto valore/sforzo per me.'}\\ \tt \\ \tt PROPOSTA: {'B91', 'A30', 'B11', 'C47', 'A47', 'A15', 'A75', 'B26', 'A19'} 
	  } 
	   } 
	   } 
	 & \\ 
 

    \theutterance \stepcounter{utterance}  

    & \multicolumn{4}{p{0.6\linewidth}}{\cellcolor[rgb]{0.8,1,0.9}{%
	\makecell[{{p{\linewidth}}}]{% 
	  \tt {\tiny [A$\rangle$GM]}  
	 ARGOMENTAZIONE: {'Sostituire C73 (347 sforzo) con A19 (862 sforzo) porta lo sforzo totale a 252+515+571+211+522+803+401+493+862 = 4630, che supera il limite di 251 unità. Questo non è fattibile. Possiamo trovare un compromesso mantenendo la struttura del set ma con aggiustamenti minori.'}\\ \tt \\ \tt RIFIUTO: {'B91', 'A30', 'B11', 'C47', 'A47', 'A15', 'A75', 'B26', 'A19'}\\ \tt \\ \tt PROPOSTA: {'B91', 'A30', 'B11', 'C47', 'A47', 'A15', 'A75', 'B26', 'C75'} 
	  } 
	   } 
	   } 
	 & & \\ 
 

    \theutterance \stepcounter{utterance}  

    & & \multicolumn{4}{p{0.6\linewidth}}{\cellcolor[rgb]{0.9,0.9,0.9}{%
	\makecell[{{p{\linewidth}}}]{% 
	  \tt {\tiny [GM$\rangle$B]}  
	 ARGOMENTAZIONE: {'Sostituire C73 (347 sforzo) con A19 (862 sforzo) porta lo sforzo totale a 252+515+571+211+522+803+401+493+862 = 4630, che supera il limite di 251 unità. Questo non è fattibile. Possiamo trovare un compromesso mantenendo la struttura del set ma con aggiustamenti minori.'}\\ \tt \\ \tt RIFIUTO: {'B91', 'A30', 'B11', 'C47', 'A47', 'A15', 'A75', 'B26', 'A19'}\\ \tt \\ \tt PROPOSTA: {'B91', 'A30', 'B11', 'C47', 'A47', 'A15', 'A75', 'B26', 'C75'} 
	  } 
	   } 
	   } 
	 & \\ 
 

    \theutterance \stepcounter{utterance}  

    & & & \multicolumn{4}{p{0.6\linewidth}}{\cellcolor[rgb]{1,0.85,0.72}{%
	\makecell[{{p{\linewidth}}}]{% 
	  \tt {\tiny [GM$\langle$B]}  
	 ARGOMENTAZIONE: {'Questa proposta ha uno sforzo totale di 252+515+571+211+522+803+401+493+339 = 4107, ben sotto il limite. Tuttavia, C75 ha un valore molto basso per me (solo 65). Preferirei includere C73 (347 sforzo, 121 valore) invece di C75 (339 sforzo, 65 valore) poiché offre un migliore rapporto valore/sforzo.'}\\ \tt \\ \tt PROPOSTA: {'B91', 'A30', 'B11', 'C47', 'A47', 'A15', 'A75', 'B26', 'C73'} 
	  } 
	   } 
	   } 
	  \\ 
 

    \theutterance \stepcounter{utterance}  

    & & \multicolumn{4}{p{0.6\linewidth}}{\cellcolor[rgb]{0.9,0.9,0.9}{%
	\makecell[{{p{\linewidth}}}]{% 
	  \tt {\tiny [A$\langle$GM]}  
	 ARGOMENTAZIONE: {'Questa proposta ha uno sforzo totale di 252+515+571+211+522+803+401+493+339 = 4107, ben sotto il limite. Tuttavia, C75 ha un valore molto basso per me (solo 65). Preferirei includere C73 (347 sforzo, 121 valore) invece di C75 (339 sforzo, 65 valore) poiché offre un migliore rapporto valore/sforzo.'}\\ \tt \\ \tt PROPOSTA: {'B91', 'A30', 'B11', 'C47', 'A47', 'A15', 'A75', 'B26', 'C73'} 
	  } 
	   } 
	   } 
	 & \\ 
 

    \theutterance \stepcounter{utterance}  

    & \multicolumn{4}{p{0.6\linewidth}}{\cellcolor[rgb]{0.8,1,0.9}{%
	\makecell[{{p{\linewidth}}}]{% 
	  \tt {\tiny [A$\rangle$GM]}  
	 ARGOMENTAZIONE: {'Questo set ha uno sforzo totale di 252+515+571+211+522+803+401+493+347 = 4115, che è ben sotto il limite e accettabile. Riconosco che C73 ha un valore leggermente migliore di C75 per te, e per me questo set include molti oggetti di alto valore come A30, B11, C47, A15 e B26. Sono disposto ad accettare questo compromesso.'}\\ \tt \\ \tt ACCORDO: {'B91', 'A30', 'B11', 'C47', 'A47', 'A15', 'A75', 'B26', 'C73'} 
	  } 
	   } 
	   } 
	 & & \\ 
 

    \theutterance \stepcounter{utterance}  

    & & & \multicolumn{2}{p{0.3\linewidth}}{\cellcolor[rgb]{0.95,0.95,0.95}{%
	\makecell[{{p{\linewidth}}}]{% 
	  \tt {\tiny [GM$|$GM]}  
	 game successful 
	  } 
	   } 
	   } 
	 & & \\ 
 

    \theutterance \stepcounter{utterance}  

    & & & \multicolumn{2}{p{0.3\linewidth}}{\cellcolor[rgb]{0.95,0.95,0.95}{%
	\makecell[{{p{\linewidth}}}]{% 
	  \tt {\tiny [GM$|$GM]}  
	 end game 
	  } 
	   } 
	   } 
	 & & \\ 
 

\end{supertabular}
}

\end{document}
